\chapter{Objetivos}

\section{Objetivo geral}

Desenvolver e analisar uma sequência didática apoiada por uma plataforma interativa para o ensino introdutório de fundamentos de \textit{machine learning}, como conteúdo de Computação, a estudantes do Ensino Médio.

\section{Objetivos específicos}

Os objetivos específicos desta pesquisa compreendem:

\begin{itemize}
    \item analisar referências sobre o ensino de inteligência artificial e \textit{machine learning} na Educação Básica, articulando o uso de representações visuais interativas à avaliação conceitual;
    \item relacionar os objetivos de aprendizagem da sequência aos eixos e às habilidades pertinentes da Computação na Educação Básica, sem supor que a norma prescreva algoritmos específicos;
    \item organizar uma sequência didática sobre regressão linear, regressão logística, \textit{k}-vizinhos mais próximos, árvore de decisão e \textit{k}-means, partindo de previsões e experimentos para observar e explicar resultados;
    \item ajustar às condições do campo de aplicação tanto os controles e explicações matemáticas da plataforma quanto os cenários sintéticos e tabelas organizadas pela turma;
    \item aplicar a sequência em uma turma de Ensino Médio em que o pesquisador é o professor;
    \item identificar o que muda ou permanece nas dificuldades dos estudantes ao explicarem como os dados e a tarefa se ligam a regras e parâmetros no resultado, confrontando o diagnóstico inicial com os registros de aula e o instrumento final;
    \item examinar como os estudantes mobilizam os conceitos estudados diante de uma situação-problema inédita em relação às aulas;
    \item sistematizar a plataforma junto à sequência didática e seus instrumentos como produto educacional.
\end{itemize}
